%%
%% This is a placeholder manuscript shell for ICSE 2028 (Research Track).
%% Refresh this manuscript before submission by adapting docs/paper/manuscript.tex
%% from the research repository
%% to the ICSE IEEE/ACM conference formatting guidelines.
%%
\documentclass[conference]{IEEEtran}

\usepackage{cite}
\usepackage{amsmath,amssymb,amsfonts}
\usepackage{algorithmic}
\usepackage{graphicx}
\usepackage{textcomp}
\usepackage{xcolor}
\usepackage{booktabs}
\usepackage{url}

\begin{document}

\title{What a Client Fetches: Default-Path Exposure of Models in Software Supply Chains}

\author{\IEEEauthorblockN{Anonymous Author(s)}
\IEEEauthorblockA{\textit{Anonymous Institution} \\
\textit{Anonymous Country} \\
anonymous@example.org}
}

\maketitle

\begin{abstract}
Software engineering increasingly relies on pre-trained machine learning models downloaded from public hubs as external components. In traditional package ecosystems like npm or PyPI, dependency managers resolve immutable packages. In contrast, model hubs rely on self-reported lineage metadata, where fine-tunes duplicate and modify weights. We examine whether derived models inherit software execution risks and licence incompatibilities from parent models. We introduce a distinction between file presence and default-path exposure: whether automated loaders execute hazardous files under default settings. In a comprehensive census of 48,247 fine-tune lineage edges spanning 16 foundational base models on Hugging Face, 0 of 16 base models exhibit default-path serialisation exposure. Across derived models, 3.17\% exhibit default-path serialisation exposure, while 54.90\% contain legacy pickle files safely ignored by modern loaders. Downstream developers introduce pickle files into safe lineages 8.5x more often than they repair them. Furthermore, 11.22\% of models declaring licences introduce incompatible terms, and 44.21\% drop licences entirely. These empirical findings demonstrate that model supply-chain defects stem predominantly from unvetted downstream repository mutations rather than inheritance from upstream foundations, providing actionable guidelines for software engineering dependency verification.
\end{abstract}

\begin{IEEEkeywords}
Software supply chain, machine learning engineering, model registries, default-path exposure, serialisation security, licence compliance
\end{IEEEkeywords}

\section{Introduction}
Modern software applications increasingly integrate machine learning models as third-party components. Rather than training models from scratch, software engineers download pre-trained weights from model registries such as the Hugging Face Hub using automated library entry points such as \texttt{AutoModel.from\_pretrained()}. 

In traditional software package ecosystems (e.g., npm, Maven, PyPI), package managers resolve dependencies via formal manifests. However, in model ecosystems, downstream developers clone and modify repositories, creating derivation trees with self-reported metadata. This paper investigates whether software security risks (arbitrary deserialisation code execution, remote execution configurations) and legal licence incompatibilities are inherited from parent foundation models or introduced downstream.

\section{Background and Threat Model}
We formalise the distinction between static file presence and default-path exposure in software supply chains. Modern loaders prioritise safe tensor formats (such as safetensors) over legacy Python pickle archives. Thus, repositories containing legacy pickle files alongside safetensors do not expose users unless specific insecure overrides are enabled.

\section{Empirical Methodology}
We conduct an empirical census of 48,247 parent--child lineage edges across 16 foundational model families. We systematically evaluate file listings, serialisation formats, remote code configuration tags, and licence compatibility matrices.

\section{Empirical Results}
Our empirical evaluation yields four primary findings:
\begin{itemize}
    \item \textbf{Absence of Parent Default-Path Exposure:} 0 of 16 base models exhibit default-path serialisation exposure; inheritance of exposure is not identified.
    \item \textbf{Presence vs Exposure Gap:} While 54.90\% ($26{,}486$ edges) contain pickle files, only 3.17\% ($1{,}528$ edges) expose clients to pickle loading by default.
    \item \textbf{Maintenance and Introduction Imbalance:} Downstream developers introduce pickle files into safe lineages 8.5x more frequently than they remediate existing pickle files.
    \item \textbf{Licence Laundering:} 11.22\% of declared child models violate upstream licence constraints, while 44.21\% omit licence documentation completely.
\end{itemize}

\section{Recommendations for Software Engineering}
We outline automated verification gates, software bill of materials (AIBOM) practices, and loader sanity checks for software engineering teams integrating external model dependencies.

\section{Conclusion}
Model supply chain defects are largely introduced by downstream maintainers rather than inherited from upstream base models. Rigorous package verification at the dependency boundary is essential.

\bibliographystyle{IEEEtran}
\bibliography{references}

\end{document}
